\section{Evaluation}
\label{sec:evaluation}

The evaluation follows the paper's fixed-budget claim through three questions.
First, within each harness's conserved budget contract, can RocksDB background
service do more useful work when placed on tenants currently creating
compaction debt? Second, is the effect online host-level placement rather than
extra capacity, static skew, local autotuning, a stale backlog signal, or a
lower-layer throttle? Third, what workload coverage, collateral cost, and
trace-qualification boundary remain after the mechanism is established?
Unless stated otherwise, percentages are online SAKI versus equal static
allocation; negative latency or bytes/write values are improvements, and
positive throughput values are improvements. In metric labels, ``High''
refers to the workload's high-pressure tenants; high-set overlap measures how
often those tenants receive SAKI's high-budget assignment. All reported
RocksDB and cgroup runs completed with zero failed tenants.

\subsection{Fixed-Budget Placement Improves High-Pressure Tenants}

The epoch-level harness isolates placement by changing RocksDB settings only
between epochs. Across five repeated 16-tenant drift trials, SAKI reduces
high-pressure write P99 by 11.7\%, raises high-pressure throughput by
38.1\%, and reduces compaction write bytes by 21.0\%
(Table~\ref{tab:evaluation-epoch-main}). Total throughput remains roughly
neutral: the mean is -1.7\%, and its CI crosses zero. All five trials improve
high-pressure write P99, high-pressure throughput, and compaction bytes; mean
high-set overlap is 3.14/4, so the policy places high budget near the moving
high-pressure set without using oracle labels.

\begin{table}[!t]
  \centering
  \caption{Epoch-level main result, \epochdrift (online vs.\ static).}
  \label{tab:evaluation-epoch-main}
  \small
  \begin{tabular}{lccc}
    \toprule
    Metric & Mean & Range & 95\% CI \\
    \midrule
    High write P99 & -11.7\% & [-17.2,-6.2] & [-18.0,-5.5] \\
    High throughput & +38.1\% & [+34.2,+42.1] & [+34.4,+41.7] \\
    Total throughput & -1.7\% & [-4.7,+2.1] & [-4.9,+1.4] \\
    Compact bytes & -21.0\% & [-24.3,-18.4] & [-23.6,-18.4] \\
    Mean high overlap & 3.14/4 & [3.14,3.14] & -- \\
    \bottomrule
  \end{tabular}
\end{table}

The continuous harness tests the stricter deployment boundary: tenants remain
alive, workload phases drift every 20 seconds, and SAKI changes only the
runtime rate-limiter budget through \ratecall. Table~\ref{tab:evaluation-continuous-main}
and Figure~\ref{fig:continuous-runtime} summarize five
\runtimedrift trials. SAKI reduces high-pressure write P99 by
27.0\%, P999 by 47.1\%, and raises high-pressure write throughput by 26.1\%.
Because SAKI completes more foreground writes, normalized bytes/write is the
main continuous efficiency metric; it falls by 11.9\%, while absolute compact
bytes change only -1.9\%. The continuous result shows the placement effect is
not an epoch artifact: the same conserved-budget loop works while tenants stay
alive and only the RocksDB rate-limiter budget changes.

\begin{table}[!t]
  \centering
  \caption{Continuous main result, \runtimedrift (five trials, online vs.\
  static).  CI half-widths in pp.}
  \label{tab:evaluation-continuous-main}
  \small
  \setlength{\tabcolsep}{3.2pt}
  \begin{tabular}{lcccc}
    \toprule
    Metric & Mean & Min & Max & 95\% CI \\
    \midrule
    High write P99 & -27.0\% & -39.2\% & -12.0\% & $\pm$12.5pp \\
    High write P999 & -47.1\% & -54.0\% & -39.3\% & $\pm$8.8pp \\
    High write throughput & +26.1\% & +21.5\% & +28.2\% & $\pm$3.5pp \\
    Total throughput & +4.3\% & +1.3\% & +6.8\% & $\pm$3.2pp \\
    Compact bytes & -1.9\% & -5.8\% & -0.1\% & $\pm$2.9pp \\
    Bytes/write & -11.9\% & -13.5\% & -9.1\% & $\pm$2.3pp \\
    Mean high overlap & 3.00/4 & 3.00 & 3.00 & -- \\
    \bottomrule
  \end{tabular}
\end{table}

\begin{figure}[t]
  \centering
  \includegraphics[width=\columnwidth]{fig3_continuous_summary.pdf}
  \caption{Continuous repeated aggregate. Each bar is normalized so positive
  means improvement over equal static; error bars show 95\% CI half-widths.}
  \Description{Horizontal bar chart of continuous-run improvements over
  static: high-pressure write P99, P999, high-pressure throughput, total
  throughput, and
  bytes/write.}
  \label{fig:continuous-runtime}
\end{figure}

\subsection{Mechanism Checks and Alternative Controls}
\label{sec:cgroup-baseline}

A fixed-budget improvement could be an artifact of a weak baseline, a stale
or fragile score instance, a non-maintenance throttle, or a lower-layer control
that already does the work. The checks in this subsection separate those
possibilities from online cross-tenant placement.

\begin{table*}[!t]
  \centering
  \caption{Mechanism checks against stronger same-budget baselines.  Rows use
  the comparison named at left, not all rows relative to static.}
  \label{tab:evaluation-stronger-baselines}
  \small
  \setlength{\tabcolsep}{4.4pt}
  \begin{tabular}{lcccccl}
    \toprule
    Comparison & High P99 & High tput & Total tput & Compact bytes & Mid tput & Overlap \\
    \midrule
    SAKI vs static & -11.7\% & +38.1\% & -1.7\% & -21.0\% & -9.1\% & 3.14/4 \\
    static\_biased vs static & -6.1\% & +17.4\% & -4.2\% & -9.9\% & -10.0\% & 1.14/4 \\
    static\_autotuned vs static & +28.8\% & -19.7\% & -1.8\% & -0.2\% & -17.3\% & 0.00/4 \\
    SAKI vs static\_biased & -6.0\% & +17.7\% & +2.8\% & -12.3\% & +1.1\% & 3.14 vs 1.14 \\
    SAKI vs static\_autotuned & -31.5\% & +71.9\% & +0.3\% & -20.8\% & +10.0\% & 3.14 vs 0.00 \\
    \bottomrule
  \end{tabular}
\end{table*}

Could equal static allocation simply be too weak a baseline?
Table~\ref{tab:evaluation-stronger-baselines} compares against a frozen skewed
allocation that helps the initial phase. Once the high-pressure set drifts,
however, its overlap falls to 1.14/4; SAKI recovers high-pressure throughput
(+17.7\% versus frozen skew), reduces compact bytes (-12.3\%), and restores
overlap to 3.14/4.

RocksDB's local auto-tuned rate limiter addresses a different problem: each
instance adapts within its configured ceiling, but unused budget is not moved
across tenants. Under the same per-tenant ceilings, \staticautotuned worsens
high-pressure P99 by 28.8\%, lowers high-pressure throughput by 19.7\%, and
has 0.00/4 high-set overlap. The same conclusion holds on the promoted
higher-pressure sensitivity point: \staticautotuned again worsens
high-pressure P99 by 29.5\% and lowers high-pressure throughput by 16.2\%
relative to static, while SAKI improves over it by 32.7\% on P99 and 59.9\%
on throughput.

A score instance that merely chases stale debt would lag the moving
high-pressure set. Pressure mode is a backlog-only online instance: it improves
high-pressure P99 ($-14.1\%$) but lags drift (overlap 1.86/4, max lag
5 windows). Demand mode adds the previous-window offered write signal, raises
overlap to 3.00/4, eliminates the 5-window lag, and raises high-pressure
throughput from +9.8\% to +26.1\%. Hybrid mode gives larger tail reductions
but still lags (overlap 2.29/4, max lag 2). The ablation matches the design
prediction in Section~\ref{sec:analysis}: backlog symptoms help, but offered
demand is what makes the rank-and-assign loop track drift rather than stale
debt.

\begin{figure*}[!t]
  \centering
  \includegraphics[width=0.88\textwidth]{fig4_score_ablation.pdf}
  \caption{Score-instance ablation. Demand mode gives strong high-pressure
  benefit while satisfying adaptation checks; pressure and hybrid instances
  improve some metrics but lag the moving high-pressure set.}
  \Description{Bar charts comparing demand, pressure, and hybrid score instances
  on normalized performance improvements, high-set overlap, and maximum lag.}
  \label{fig:score-ablation}
\end{figure*}

Is the runtime actuator really a maintenance I/O knob? The four-tenant
coverage microstudy in Section~\ref{sec:implementation} shows measured
limiter traffic rising monotonically with the configured budget and
compact-plus-flush output tracking it (coverage ratio 0.99--1.02) until the
workload's maintenance demand saturates near 7.3\,MB/s. This confirms the
continuous actuation path used by SAKI.

Is the result tied to one fragile score instance?
Section~\ref{sec:saki-design} treats the score
$s_i = 6.00\,D/1000 + 0.60\,G/1000 + 0.05\,T + 0.02\,B + 0.02\,P + 0.30\,L$
as an auditable reference instance, not an optimum. First, we replay the 35
post-warmup control decisions from the five continuous main online traces using
the previous-window observations that the controller saw. The replay exactly
reproduces the recorded online high-budget set in all 35 decisions. The demand
term alone also reproduces the reference top-four set in all 35 decisions;
anchor-dominant perturbations keep exact or near-exact top-four agreement,
while a pressure-only negative control matches only 10/35 decisions
(Table~\ref{tab:ranking-margin}). The reference score's top-four boundary
margin is positive in every replayed decision (mean 5.30, min 0.39), so the
main assignment is a property of a demand-aware score family rather than an
artifact of a single coefficient vector.

\begin{table}[!t]
  \centering
  \caption{Ranking-margin replay on the continuous main online traces
  (35 post-warmup decisions). Each score instance is recomputed from the
  previous-window observations; Jaccard compares its top-four set with the
  reference demand-mode top-four set.}
  \label{tab:ranking-margin}
  \small
  \setlength{\tabcolsep}{3.2pt}
  \begin{tabular}{lccc}
    \toprule
    Score instance & Exact & Jaccard & True overlap \\
    \midrule
    demand only & 35/35 & 1.00 [1.00,1.00] & 3.00/4 \\
    anchor $\times 0.5$ & 31/35 & 0.95 [0.33,1.00] & 2.86/4 \\
    anchor $\times 2$ & 35/35 & 1.00 [1.00,1.00] & 3.00/4 \\
    residual $\times 0.5$ & 35/35 & 1.00 [1.00,1.00] & 3.00/4 \\
    residual $\times 2$ & 31/35 & 0.95 [0.33,1.00] & 2.86/4 \\
    pressure-only & 10/35 & 0.61 [0.14,1.00] & 2.03/4 \\
    \bottomrule
  \end{tabular}
\end{table}

Second, we rerun nearby score instances on \runtimedrift by perturbing anchor
weight ($\times 0.5$, $\times 2$) and residual scale
($\times 0.5$, $\times 2$, $\times 0$, the last being anchor-only), with three
seeds each paired with the published paper-static trials a/b/c.
Table~\ref{tab:coefficient-ablation} shows that all five demand-aware
perturbations keep the headline direction, maintain overlap at
2.86--3.00/4, and keep max lag to 0--1 window. The \textsc{anchor-only} row,
where the residual scale is zero, is the runtime demand-only case: it confirms
the replay result that offered demand alone fixes the high-budget set, with
overlap 3.00/4 and zero lag. The residuals matter in the lower bound:
anchor-only high-throughput improvement drops as low as \(+12.6\%\), whereas
restoring half-strength residuals raises the minimum to \(+21.3\%\). This is consistent
with the division of labor in Section~\ref{sec:saki-design}: the demand anchor
tracks the active high-pressure set, while bounded pressure residuals appear to
stabilize near-boundary tenants whose maintenance state has not fully drained.
This is the empirical counterpart of Proposition~1: nearby demand-aware
instances keep the same placement behavior, whereas removing demand collapses
into the pressure-only case (Figure~\ref{fig:score-ablation}). Because the
ablation rows pair with only three static trials, row-to-row differences should
be read as consistency checks rather than tuning claims.

\begin{table*}[!t]
  \centering
  \caption{Score-family robustness around the demand-mode SAKI instance on
  \runtimedrift ($n{=}3$ per row, paired with paper-static a/b/c). Each row
  perturbs only the anchor or residual coefficients; all other settings are
  fixed. Cells show $\mathrm{mean}\,[\min,\max]$ percent of paired static;
  Top-4 Jaccard averages the post-warmup window-by-window high-set overlap
  against published SAKI decisions. The \textsuperscript{*} row reproduces
  the published $n{=}5$ main continuous result; its static pairing covers all
  five trials, not the three used by the ablation rows.}
  \label{tab:coefficient-ablation}
  \small
  \setlength{\tabcolsep}{4pt}
  \begin{tabular}{lccccc}
    \toprule
    Perturbation & High P99 & High tput & Overlap & Max lag & Top-4 Jaccard \\
    \midrule
    anchor $\times 0.5$ (3.00)        & -30.3\% [-38.1,-23.6] & +22.4\% [+18.9,+27.1] & 2.90/4 & 1 & 0.94 [0.88,1.00] \\
    anchor $\times 2$ (12.00)         & -27.5\% [-42.4,-18.2] & +23.7\% [+22.2,+26.2] & 3.00/4 & 0 & 1.00 [1.00,1.00] \\
    residual $\times 0.5$             & -30.0\% [-39.8,-16.7] & +24.3\% [+21.3,+28.2] & 3.00/4 & 0 & 1.00 [1.00,1.00] \\
    residual $\times 2$               & -35.8\% [-49.6,-25.7] & +25.0\% [+20.9,+27.1] & 2.86/4 & 1 & 0.94 [0.89,1.00] \\
    anchor only (residual$=0$)        & -31.9\% [-40.2,-26.2] & +21.5\% [+12.6,+26.0] & 3.00/4 & 0 & 1.00 [1.00,1.00] \\
    \midrule
    fixed SAKI (n=5, paper main)\textsuperscript{*} & -27.0\% & +26.1\% & 3.00/4 & 0 & 1.00 (self) \\
    \bottomrule
  \end{tabular}
\end{table*}

Finally, would a lower-layer control point be equivalent?
On the aligned epoch workload, online cgroup-v2 control improves high-pressure
throughput by 12.5\% relative to equal cgroup throttling, but the P99 CI
crosses zero (-4.4\%) and compaction bytes increase by 3.1\%. An oracle cgroup
assignment can express priority when the high-pressure set is known (P99
-9.1\%, throughput +15.4\%), so the layer is useful as an enforcement
envelope. It does not provide the LSM-aware budget-placement signal that SAKI
uses.

\subsection{Controlled Coverage Around the Anchor Workloads}
\label{sec:workload-matrix}

The next experiments test local coverage around the controlled anchor
workloads while preserving the same-budget contract. The epoch sensitivity
sweep changes drift speed and high-write pressure around \epochdrift. These
rows are local coverage checks: the promotion rule keeps capacity-boundary and
failed-tenant cases from obscuring the comparison, while the headline claim
rests on the repeated anchors above. All promoted points move in the headline
direction: high-pressure P99 improves by 6.9--12.9\%, high-pressure throughput
improves by at least 33.0\%, and compaction bytes fall by 17.6--21.3\%. The
faster-drift point remains CI-strict on high-pressure benefit (P99 -6.9\%,
throughput +33.9\%) with an aggregate-throughput cost (-5.0\%), matching the
expected responsiveness limit of a previous-window controller.

The continuous workload matrix changes one axis at a time around
\runtimedrift: 4\,KB values, tripled point-read QPS, faster drift at
1.5 tenants/window, and an 8-tenant cluster with two high-pressure tenants.
All four variants were promoted to five trials and retain the headline
direction: high-pressure P99 improves in every row and high-pressure
throughput improves by 20.1--30.4\%
(Table~\ref{tab:evaluation-workload-matrix}). Three rows are CI-strict wins
on both primary high-pressure metrics; \readthreex remains CI-strict on high
throughput, but its P99 confidence interval crosses zero. The matrix also
names the trade-off rather than hiding it: \valuefourk's P999 mean is negative
but not statistically separated from static, reflecting sparse
heavy-compaction bursts under a conserved rate-limit ceiling rather than a P999
regression. \readthreex is a point-read-QPS perturbation, not a scan-heavy or
read-dominant SLO result.

\begin{table*}[!t]
  \centering
  \caption{Local workload matrix, five trials per variant, online vs.\ static
  (percent changes, 95\% CI half-widths in pp).  Zero failed tenants in all
  runs.}
  \label{tab:evaluation-workload-matrix}
  \small
  \setlength{\tabcolsep}{2.8pt}
  \renewcommand{\arraystretch}{1.08}
  \begin{tabular}{lcccccr}
    \toprule
    Variant & High P99 & High P999 & High tput & Total tput & Bytes/write & Overlap \\
    \midrule
    \valuefourk   & -28.6\%$\pm$16.2 & -7.1\%$\pm$20.6  & +20.1\%$\pm$4.7  & +0.6\%$\pm$1.5  & -9.0\%$\pm$3.0  & 3.00/4 \\
    \readthreex   & -15.4\%$\pm$17.5 & -38.4\%$\pm$27.6 & +30.4\%$\pm$7.6  & -0.2\%$\pm$4.0  & -10.2\%$\pm$5.8 & 3.00/4 \\
    \fastdrift    & -17.6\%$\pm$10.1 & -1.1\%$\pm$46.6  & +20.5\%$\pm$10.3 & +4.4\%$\pm$3.4  & -3.3\%$\pm$4.4  & 2.57/4 \\
    \tenantseight & -42.2\%$\pm$12.7 & -52.6\%$\pm$16.9 & +28.7\%$\pm$5.8  & +7.0\%$\pm$1.9  & -11.3\%$\pm$3.4 & 1.57/2 \\
    \bottomrule
  \end{tabular}
\end{table*}

The window length is also a mechanism setting, not a tuned constant.
Table~\ref{tab:evaluation-window-length} repeats the continuous anchor at
10, 20, and 40 seconds with three seed pairs each, holding duration, offered
load, drift distance, and the 96\,MB/s aggregate budget fixed. The 20-second
anchor is the only point with a CI-strict high-P99 win (-26.9\%); 10 seconds
keeps high throughput positive but makes both P99 and throughput noisy and
raises LOW P99, while 40 seconds lowers overlap to 2.00/4 and leaves
adaptation lag unbounded in all three seeds. This supports using 20 seconds as
a mechanism-isolation anchor and treats the other two points as responsiveness
boundaries.

\begin{table}[!t]
  \centering
  \caption{Window-length sensitivity around the continuous anchor (three
  seed pairs per row, online vs.\ static; CI half-widths in pp).}
  \label{tab:evaluation-window-length}
  \small
  \setlength{\tabcolsep}{3.0pt}
  \begin{tabular}{lcccc}
    \toprule
    Window & High P99 & High tput & Overlap & LOW P99 \\
    \midrule
    10\,s & +6.8\%$\pm$39.2  & +23.4\%$\pm$58.1 & 3.47/4 & +110.5\%$\pm$367.8 \\
    20\,s & -26.9\%$\pm$14.0 & +16.1\%$\pm$8.7  & 3.00/4 & +24.6\%$\pm$63.6 \\
    40\,s & -14.4\%$\pm$44.8 & +14.5\%$\pm$19.8 & 2.00/4 & -36.3\%$\pm$84.2 \\
    \bottomrule
  \end{tabular}
\end{table}

\subsection{Collateral Cost and Liveness}

SAKI is a targeted reallocation policy, not a Pareto-improvement claim.
Tier-level accounting makes the redistribution explicit. In the continuous
aggregate, HIGH tenants receive P99 -27.0\%, P999 -47.1\%, and write
throughput +26.1\%; LOW tenants absorb P99 +1.4\%, P999 +2.6\%, and write
throughput -3.0\%. The epoch-level pattern is similar: HIGH throughput rises
by 38.1\%, while MID and LOW throughput change by -9.1\% and -1.2\%. Total
throughput remains close to static in both harnesses, so the benefit is a
budget-placement trade-off rather than added capacity.

\begin{figure*}[!t]
  \centering
  \includegraphics[width=0.82\textwidth]{fig8_collateral_summary.pdf}
  \caption{Tier-wise benefit and collateral cost. Cells report SAKI's
  improvement over static (positive is better), showing the intended
  high-pressure benefit and the measured redistribution cost to lower tiers.}
  \Description{Two-panel matrix for continuous and epoch runs, showing P99,
  P999, and throughput changes by high, mid, and low tiers.}
  \label{fig:collateral-summary}
\end{figure*}

Read-side and CPU measurements show that the bytes/write improvement does not
hide higher compaction CPU per completed write: CPU/write falls by 12.8\%
overall. The tier values are -1.6\% for HIGH, -8.0\% for MID, and -33.5\%
for LOW. Read-side effects are bounded in this read-light workload (overall
read P99 +7.0\%, read P999 -10.6\%, read throughput +0.5\%), but they are
collateral accounting rather than evidence for read-dominant SLOs. A separate
360-second continuous run with 16 drift events reports high P99 $-13.9\%$,
high P999 $-19.6\%$, high write throughput $+11.6\%$, bytes/write $-3.4\%$,
overlap $3.00/4$, max lag $1$~window, and zero tenant failures---a liveness
and tracking sanity check at the longer horizon.

\subsection{Trace Qualification Boundary}
\label{sec:public-trace}

The controlled experiments above establish the fixed-budget placement
mechanism. Public traces become meaningful controller-validation workloads
when they preserve tenant identity and drift, expose actuatable supply under
the fixed budget, and create enough RocksDB maintenance stress to exercise
compaction-budget placement. We audited two public trace families under the
same 11/7/3\,MB/s, 112\,MB/s budget contract and use them to identify which of
those conditions are missing.

Meta CacheLib \texttt{kvcache/202401}~\cite{cachelib_trace} is useful as a
trace-shaped replay and tracking smoke, but it records cache-front-end
requests rather than RocksDB compaction events. Five selected drifting
16-tenant segments completed with zero failed tenants and preserved drift, yet
the replay failed the engine-stress gate: offered writes were 110--143\,MB/s,
completed logical writes were only 6.7--7.3\,MB/s, compaction output plateaued
near 29--31\,MB/s, and no segment entered the stress-pass subset. Meta
Tectonic/Baleen storage 0.1\%~\cite{baleen_fast24} exposes the complementary
condition. Aggregate scaled write supply initially looked sufficient, but most
windows have only one to three active write tenants. Under the static all-mid
policy, the actuatable supply \(\sum_i \min(7,\mathrm{traceWrite}_i)\) is far
below the aggregate trace rate; across 146 independent candidates, zero have
six or more windows above 70\,MB/s or mean static-actuatable supply above
70\,MB/s. These audits are therefore trace-qualification results rather than
controller comparisons.
